% Source: https://artofproblemsolving.com/community/c3h2779451_2017_mbmt_team_round__montgomery_blair_math_tournament
% Stable identity: aops-topic-2779451-q115
For $n=\prod p_i^{e_i}$, divisor independence gives $\sum_{d\mid n}\tau(d)=\prod_i(1+2+\cdots+(e_i+1))=\prod_i\dfrac{(e_i+1)(e_i+2)}{2}$. The only triangular factors greater than 1 dividing 108 are 3,6,36, corresponding to exponents 1,2,7. Products equaling 108 are $36\cdot3$ and $6\cdot6\cdot3$. Thus $n$ has form $p^7q$ or $p^2q^2r$, with distinct primes. The first type is at most 896: $p$ must be 2, and $q\le7$. For the second type, $p<q$ implies $pq\le\sqrt{\dfrac{999}{2}}<23$, so the possible squared-prime pairs are $(2,3),(2,5),(2,7),(2,11),(3,5),(3,7)$. Maximizing the distinct prime $r$ gives respectively 828,700,980, no admissible $r$,450,882. Hence $\boxed{980=2^2\cdot7^2\cdot5}$.
